\section{Onde aprofundar}

As fontes principais deste capítulo, na ordem em que aparecem. Acesso
(aberto ou pago) e links também na bibliografia.

\subsection*{Especificações}
\begin{itemize}
\obra{livro:gosling2000-jls2}{A JLS 2ª ed., de onde vêm as regras de
  threads, locks e memória antes da JSR~133 (cap.~17).}
\obra{livro:lindholm1999}{A JVMS 2ª ed., com os quatro passes do
  verificador (seção~\ref{sec:verificacao}).}
\obra{doc:jls25}{A JLS atual, para comparar o cap.~17 com o da 2ª ed.}
\obra{doc:jvms25}{A JVMS atual, com o verificador por type checking e as
  flags de campo e de método mostradas nas listagens.}
\end{itemize}

\subsection*{Decisões e história}
\begin{itemize}
\obra*{paper:java-language-environment}{O white paper que lista os
  objetivos de projeto (seção~\ref{sec:objetivos-1995}).}
\obra*{paper:hoare1974-monitors}{O artigo de monitores; o white paper
  diz que o paradigma de monitor do Java vem de Hoare.}
\obra*{paper:brinch-hansen1999-insecure}{A crítica de 1999 aos monitores
  do Java (seção~\ref{sec:monitores}).}
\obra*{paper:manson-tese-jmm-2004}{A tese, aberta, que explica o que
  estava errado no modelo de memória antigo e o modelo novo.}
\obra*{doc:jsr133-faq}{Os sintomas do modelo antigo e o
  \texttt{volatile} e o \texttt{final} do novo, em perguntas e respostas.}
\obra*{paper:lea2000-forkjoin}{O framework Fork/Join e o work-stealing,
  antes de entrarem no JDK (seção~\ref{sec:juc}).}
\obra*{artigo:kegel-c10k}{O texto que pergunta como atender dez mil
  clientes ao mesmo tempo (seção~\ref{sec:muitos-requests}).}
\obra*{doc:jep444}{Por que as virtual threads existem: dar escala, não
  velocidade (seção~\ref{sec:threads-ate-vt}).}
\end{itemize}

\subsection*{Para o código do dia a dia}
\begin{itemize}
\obra*{livro:goetz2006}{Visibilidade, publicação segura e
  \texttt{java.util.concurrent} na prática.}
\obra{doc:javase25-juc}{O javadoc do pacote, com as coleções
  concorrentes e quando preferi-las a um mapa sincronizado.}
\end{itemize}
